\documentclass[11pt]{article}
\usepackage[margin=25mm]{geometry}
\usepackage{fontspec}
\setmainfont{Latin Modern Roman}
\usepackage{amsmath,amssymb,mathtools}
\usepackage{unicode-math}
\setmathfont{Latin Modern Math}
\usepackage{xurl,hyperref}
\hypersetup{colorlinks=true,urlcolor=blue}
\setcounter{secnumdepth}{0}
\setlength{\parindent}{0pt}
\setlength{\parskip}{5pt}
\setlength{\emergencystretch}{3em}
\newcommand{\sref}[2]{\hyperlink{source:#1:#2}{[S#2]}}
\newcommand{\cataloguescope}[1]{\paragraph{Recorded target.}#1}
\newcommand{\R}{\mathbb R}
\newcommand{\C}{\mathbb C}
\newcommand{\N}{\mathbb N}
\newcommand{\Z}{\mathbb Z}
\newcommand{\Q}{\mathbb Q}
\newcommand{\F}{\mathbb F}
\newcommand{\eps}{\varepsilon}
\begin{document}
% UnsolvedMath ID 30000361; source catalogue code OWR-1114-002
\setcounter{section}{30000360}
\section{Heterotic Mirror Symmetry for Bundled Calabi–Yau Manifolds}\label{30000361}
{\small UnsolvedMath ID 30000361 · Mathematical Physics}
\subsection{Definitions and mathematical statement}

\paragraph{Shared notation.}$\mathbb N=\{1,2,\ldots\}$; the standard number fields have their usual meanings. All additional conventions are stated locally.


\paragraph{Definitions and conventions.} A heterotic $(0,2)$ compactification pairs a Calabi--Yau manifold $X$ with a holomorphic vector bundle $E$ carrying the left-moving gauge data, subject to the physical consistency conditions of the model. A $(0,2)$ mirror pair is expected to exchange the combined complex, Kähler, and bundle moduli and corresponding sheaf-cohomology data, and to reduce to ordinary mirror symmetry when $E\simeq TX$. A Fourier--Mukai transform is an integral functor between derived categories with kernel on the product of two varieties; an elliptic K3 surface is a K3 surface fibred by elliptic curves.

\paragraph{Source scope and qualifications.} The report asks for a general meaningful formulation as well as a procedure. It does not define a unique universal mirror functor or state that all bundles meet the string anomaly and stability conditions. The elliptic-K3 proposal is a particular included question.

\paragraph{Catalogue mathematical target.}
\paragraph{Catalogue wording retained for traceability.}
Construct a general formulation and construction procedure for heterotic $(0,2)$ mirror pairs of Calabi–Yau manifolds with holomorphic vector bundles; in particular, determine whether such mirror symmetry can be realized functorially by Fourier–Mukai transforms in the elliptically fibered K3 setting.

\paragraph{Saved literature context.} The supplied catalogue's assessment dated 2026-08-21 records: Fourier--Mukai descriptions of special K3 mirror situations and treatments of selected gauge bundles exist, but the general (0,2) mirror formulation and the proposed construction for arbitrary structure groups remain open. The assessment is retained as dated provenance; it is not a new comprehensive literature or open-status certification.
\subsection{Short English statement}
Formulate and construct heterotic mirror pairs with vector bundles, including the proposed Fourier--Mukai description for elliptic K3 surfaces.
\subsection{Research attempt}
\paragraph{Model, date and investigated scope.}
OpenAI GPT-6 Astra (configured model \texttt{gpt-6-astra}), ultra reasoning; generated 11 October 2026. Computes anomaly and slope-stability constraints for the split rank-two bundle ansatz, showing why topological matching alone is inadequate.

\paragraph{Sources and literature check.}
The complete supplied definition, including source qualifications and the dated literature assessment, was reviewed. The calculation below is an elementary restricted attempt motivated by [S2]; it does not claim novelty or a fresh literature-status determination.

\paragraph{Context and investigated scope.} The complete supplied definition, including source qualifications and the dated literature assessment, was reviewed. The calculation below is an elementary restricted attempt motivated by [S2]; it does not claim novelty or a fresh literature-status determination.

\paragraph{A consistency test for split-bundle constructions.} Restrict to a complex surface $X$ with $c_1(TX)=0$ and a candidate rank-two bundle $E=L\oplus L^{-1}$. Put $\ell=c_1(L)$. Multiplication of total Chern classes gives
\[
 c(E)=(1+\ell)(1-\ell)=1-\ell^2,\quad c_1(E)=0,\quad c_2(E)=-\ell^2.
\]
In the no-extra-source anomaly model imposing $c_2(E)=c_2(TX)$, a necessary condition is therefore $-\ell^2=c_2(TX)$. It is an imposed physical consistency condition in this restricted model, not an assertion that every heterotic compactification has this exact equation.

\paragraph{Stability obstructs the easiest ansatz.} For a Kähler class $\omega$, define slope of a rank-$r$ bundle by $\mu_\omega(F)=r^{-1}\int_Xc_1(F)\wedge\omega$. The two line summands have opposite slopes $s$ and $-s$, whereas $\mu_\omega(E)=0$. Strict slope stability requires every proper subbundle to have slope below zero. If $s\ne0$, one summand violates this; if $s=0$, both give equality and strict stability still fails. Thus this elementary split ansatz never produces a strictly stable rank-two bundle, although at zero slope it may satisfy a weaker polystability requirement.

\paragraph{Why the attempted shortcut stops.} Chern-number matching alone cannot turn a candidate Fourier--Mukai construction into a physical mirror: even this simple family separates topological consistency from stability. One must specify allowed stability, ensure the transformed object is an appropriate bundle or physical complex, and compare moduli and spectra. No functor implementing those requirements on elliptic K3 surfaces is constructed, and the general $(0,2)$ formulation, including reduction to the tangent-bundle case, remains beyond the calculation.

\paragraph{Verification and adversarial audit.}
Checked Chern-class signs, both signs of slope and the zero-slope boundary; distinguished stability, polystability and the chosen anomaly model. Model self-audit only; no independent human proof checking. Prompts A and B were applied to this canonical record and its qualified source scope. Checked Chern-class signs, both signs of slope and the zero-slope boundary; distinguished stability, polystability and the chosen anomaly model. Retained unresolved status and zero original-target progress; No Fourier--Mukai transform preserving the required physical conditions and exchanging the full moduli/spectrum data has been constructed.

Coordinator checked the split-bundle Chern classes and both slope-sign cases. The calculation obstructs strict stability of that ansatz only. This additional check was model-assisted, not independent human verification.

\paragraph{Tools, process and contributions.}
Python 3 was used for file preparation and structural checks; no numerical experiment or theorem prover was used. Contribution prompt A was followed by prompt B. GPT-6 Astra with ultra reasoning drafted the research and performed its stated self-audit; any additional coordinator contribution is identified in the verification and audit. A GPT-6-based Codex coordinator assembled the manuscript and checked repository integration; its exact variant and reasoning configuration were not exposed. No human mathematical contribution, independent human proof check or formal proof certification is supplied. Complete derivations and the consulted source references are included here; private reasoning transcripts are not published.

\paragraph{Final status and first remaining gap.}
\textbf{UNRESOLVED}. No Fourier--Mukai transform preserving the required physical conditions and exchanging the full moduli/spectrum data has been constructed.

\section{Completion Estimate}
\noindent\textbf{COMPLETION: 0\%}
The calculation is an auxiliary or restricted attempt and establishes no newly verified portion of the full source target. No Fourier--Mukai transform preserving the required physical conditions and exchanging the full moduli/spectrum data has been constructed.

\subsection{Sources}
\begin{itemize}
\item[\hypertarget{source:30000361:1}{S1}] ULAM.AI and the UnsolvedMath Contributors, UnsolvedMath, record 30000361 (OWR-1114-002). Catalogue curation under CC BY 4.0; definitions and English exposition added locally.
\url{https://huggingface.co/datasets/ulamai/UnsolvedMath}.
\item[\hypertarget{source:30000361:2}{S2}] Mini-Workshop: Heterotic Strings, Derived Categories, and Stacks (2005), Oberwolfach Reports, DOI 10.4171/owr/2005/53, pp. 3019--3060; official report source and mathematical context.
\url{https://doi.org/10.4171/owr/2005/53}.
\item[\hypertarget{source:30000361:3}{S3}] Official EMS Press PDF of the same Oberwolfach report; PDF pages 9--10, printed pages 3027--3028.
\url{https://ems.press/content/serial-article-files/46027}.
\item[\hypertarget{source:30000361:4}{S4}] string compactifications on stacks; reference in the saved dated catalogue literature review.
\url{https://arxiv.org/abs/hep-th/0502053}.
\end{itemize}
\label{end:30000361}
\end{document}
